% Generated draft: requires mathematical and source review.
\documentclass[11pt]{article}
\usepackage[margin=25mm]{geometry}
\usepackage{fontspec}
\setmainfont{Latin Modern Roman}
\usepackage{amsmath,amssymb,mathtools,unicode-math}
\setmathfont{Latin Modern Math}
\usepackage{enumitem,needspace,xurl,hyperref}
\setlength{\parindent}{0pt}
\setlength{\parskip}{5pt}
\setlength{\emergencystretch}{3em}
\newcommand{\N}{\mathbb{N}}
\newcommand{\Z}{\mathbb{Z}}
\newcommand{\Q}{\mathbb{Q}}
\newcommand{\R}{\mathbb{R}}
\newcommand{\C}{\mathbb{C}}
\newcommand{\F}{\mathbb{F}}
\newcommand{\A}{\mathbb{A}}
\newcommand{\PP}{\mathbb{P}}
\newcommand{\E}{\mathbb{E}}
\newcommand{\eps}{\varepsilon}
\newcommand{\dd}{\,\mathrm{d}}
\begin{document}
\section*{UnsolvedMath ID: 20002085}
\textbf{Source problem code:} catalogue entry
\subsection{Definitions and mathematical statement}
Question 1.27 records Hutchings's suggestion to extend methods from his talk to compute embedded contact homology of \(S^1\times S^2\), torus bundles over \(S^1\), and unit cotangent bundles of surfaces of genus greater than one. It also asks to understand holomorphic curves in \(\mathbb R\times T^3\) using amoebas or related descriptions. The talk's methods and the contact structures, homology classes and chain conventions of the requested computations are not supplied.

The paired aid treats a curve-level special case. Let \(K\subset\mathbb R^2\) be a smooth strictly convex body with the origin in its interior, and parametrize its boundary positively as \((a(\theta),b(\theta))\). On \(T^3\), with \(x,y\) of period one, use
\[
 \lambda_K=a(\theta)dx+b(\theta)dy,\qquad
 h_K(v)=\max_{p\in K}\langle p,v\rangle.
\]
The interior-origin condition ensures the nonvanishing needed for this contact form; mere containment on the boundary would not suffice. A closed Reeb orbit lies at fixed \(\theta\) and has nonzero lattice displacement \(v\in\mathbb Z^2\), with multiplicity included in \(v\). The reported action formula is
\[
 \mathcal A(\gamma)=\int_\gamma\lambda_K=h_K(v).
\]
Its normal direction selects the supporting point of \(K\), and support-function homogeneity accounts for covers.

For a finite-energy holomorphic current adapted to this symplectization and asymptotic to closed Reeb orbits, let positive-end vectors be \(v_i\) and negative-end vectors \(w_j\), including multiplicities. The aid reports
\[
 \sum_i v_i=\sum_j w_j,
\]
\[
 0\leq E_{d\lambda}(C)=\int_Cd\lambda_K
       =\sum_i h_K(v_i)-\sum_jh_K(w_j).
\]
Thus the support-perimeter functional, defined here as the corresponding sum of support-function values, decreases from positive to negative orbit sets.

If there is exactly one positive end with vector \(v\), conservation gives \(v=\sum_jw_j\). Subadditivity yields
\[
 h_K(v)\leq\sum_jh_K(w_j),
\]
so the nonnegative energy is zero. The reported conclusion is that the current is supported on trivial Reeb cylinders. This permits covered or branched-cover maps with that image; it does not claim every parametrized curve is an unbranched cylinder.

This obstruction does not compute an ECH differential. Orbit-set generators can have several positive ends, and an ECH chain model requires the relevant nondegeneracy/perturbation and counting setup. The toric contact forms naturally have orbit families, so conclusions for the unperturbed form cannot simply be transferred to every perturbation. No ECH computation for the other listed manifolds or complete amoeba correspondence is supplied by this special case.
\subsection{Short English statement}
The source proposes several ECH computations and an amoeba description of curves in \(\mathbb R\times T^3\). The supplied aid proves a narrower convex-toric obstruction: action is a support function, total lattice displacement is conserved, and a one-positive-end current must have zero transverse energy and trivial-cylinder support. This is not an ECH differential calculation.
\subsection{Sources}
\begin{itemize}
\item Local source record: catalogue entry, catalogue entry; canonical \texttt{aim-geometry-notes.json}, \texttt{notes[422]}, Question 1.27. Title: ``A convex support-function obstruction for holomorphic curves in R x T\^{}3.''
\item AIM workshop notes \emph{Holomorphic curves in contact geometry}, supplied reference \url{https://aimath.org/WWN/contactgeom2/contactgeom2.pdf}. The question is attributed to Hutchings's talk and appears under ``Fourth day'' according to the supplied verification. This draft does not access a PDF or HTML source.
\item Research attempt 1: convex support-function action formula, lattice conservation, perimeter inequality and one-positive-end obstruction. No full authoritative external bibliographic citation is supplied.
\item Registry dated 14 August 2026: ``partially\_solved,'' ``PARTIAL-PROGRESS,'' statement status ``exact''; attempt partial, rigor ``proved,'' novelty confidence low. The machine-generated, non-peer-reviewed aid requires independent verification before citation.
\end{itemize}
\end{document}
